\section{When the full Frank--Wolfe gap already bounds the error}
\label{gapregime:section}
The outer maximization in Theorem~\ref{boundarygap:main} is unnecessary
in an explicit range of observed gaps.  The range is exact for the
stated function class.  Throughout this section, $\mathcal B$ denotes
the full-gap envelope in \eqref{boundarygap:definition}, with the
same shifted-RSI and norm-Hessian assumptions.

For $H_1>0$ and $g\ge0$, put $b=H_1g$ and let $\tau_g\in(1/2,1]$
be the unique solution of
\begin{equation}\label{gapregime:switch}
 b\tau_g^2+2H_0\tau_g=H_0+m+b/2.
\end{equation}
An expression that also applies at $b=0$ is
\begin{equation}\label{gapregime:tau}
 \tau_g=
 \frac{H_0+m+b/2}{H_0+\sqrt{H_0^2+b(H_0+m)+b^2/2}}.
\end{equation}
Define the algebraic quantity
\begin{equation}\label{gapregime:coefficient}
 \mathcal S(g)=H_0\left(\frac12-\tau_g^2\right)
       +\frac{H_1g}{3}\left(1-2\tau_g^3\right).
\end{equation}

\begin{theorem}[The exact range of the ordinary gap certificate]
\label{gapregime:main}
Let $m>0$, $H_0\ge m$, and $H_1>0$.  For every $g>0$,
\begin{equation}\label{gapregime:equivalence}
 \boxed{\quad \mathcal B(g)=g
              \quad\Longleftrightarrow\quad \mathcal S(g)\le0.\quad}
\end{equation}
When $\mathcal S(g)\le0$, every positive radius in the defining
maximization satisfies $V(g,\rho)<g$.  Thus the envelope equals the
limit at $\rho=0$; no positive radius attains it.  In particular,
every admissible nonoptimal source with full gap $g$ satisfies
$\delta(x)<g$ in this range.

The function $\mathcal S$ is continuously differentiable and strictly
increasing on $[0,\infty)$, with
\begin{align}
 \mathcal S(0)
   &=\frac{H_0^2-2mH_0-m^2}{4H_0},
       \label{gapregime:initial}\\
 \mathcal S'(g)
   &=H_1\left(\frac{\tau_g^3}{3}
                    -\frac{\tau_g}{2}+\frac13\right)>0.
       \label{gapregime:derivative}
\end{align}
Consequently, if $m\le H_0<(1+\sqrt2)m$, there is a unique
$g_{\rm c}>0$ with $\mathcal S(g_{\rm c})=0$, and
\[
 \mathcal B(g)=g\quad(0\le g\le g_{\rm c}),\qquad
 \mathcal B(g)>g\quad(g>g_{\rm c}).
\]
If $H_0\ge(1+\sqrt2)m$, then $\mathcal B(g)>g$ for every $g>0$.
For $H_0=m$ the threshold has the exact form
\begin{equation}\label{gapregime:unit-threshold}
 \boxed{g_{\rm c}=(2+\sqrt6)\,m/H_1.}
\end{equation}
The proof is in Appendix~\ref{gapregime:proof}.
\end{theorem}

\begin{corollary}[Radius monotonicity and an algebraic upper bound]
\label{gapregime:radius-monotone}
Under the theorem's assumptions, fix $g>0$ with $\mathcal S(g)\le0$.
Then $\rho\mapsto V(g,\rho)$ is strictly decreasing on $[0,g/m]$.
More precisely, for $0<\rho_1<\rho_2\le g/m$,
\begin{equation}\label{gapregime:deficit-ratio}
 \frac{g-V(g,\rho_1)}{\rho_1}
       \le\frac{g-V(g,\rho_2)}{\rho_2}.
\end{equation}
Consequently,
\begin{equation}\label{gapregime:tangent-upper}
 V(g,\rho)\le g+\rho\mathcal S(g),
       \qquad 0\le\rho\le g/m.
\end{equation}
On any restricted radius interval $[\rho_-,\rho_+]\subseteq[0,g/m]$,
the maximum is therefore the left endpoint value $V(g,\rho_-)$.
It can be evaluated with at most one switch root, with no root needed
when $\rho_-=0$, or bounded algebraically by $g+\rho_-\mathcal S(g)$.
\end{corollary}

For general $m\le H_0<(1+\sqrt2)m$, set
$d=-\mathcal S(0)>0$.  The threshold obeys the explicit bounds
\begin{equation}\label{gapregime:threshold-bracket}
 \frac{6d}{H_1}\le g_{\rm c}
 \le\frac{3d}{H_1(1-1/\sqrt2)}.
\end{equation}
There is no need to find $g_{\rm c}$ to test a source: evaluate
\eqref{gapregime:tau}--\eqref{gapregime:coefficient} at its observed
gap.  The test uses a square root and arithmetic, with no switch
root search and no radius grid.  When $H_1=0$, the already proved
formula \eqref{boundarygap:constant} gives the corresponding
gap-only range: $\mathcal B(g)=g$ for all $g\ge0$ exactly when
$H_0\le(1+\sqrt2)m$.

\paragraph{A supplied diameter leaves the threshold unchanged.}
For every fixed $D>0$, the same equivalence holds with
$\mathcal B_D(g)$ in place of $\mathcal B(g)$.
The nonpositive case follows from
$g\le\mathcal B_D(g)\le\mathcal B(g)$.
In the positive case, \eqref{gapregime:variation} exceeds $g$ at
arbitrarily small positive radii, which also satisfy $\rho\le D^2$.

\paragraph{Certified source observations.}
If a paid source query gives $0\le G(x)\le B$ and
$\mathcal S(B)\le0$, then
\begin{equation}\label{gapregime:inexact}
 \delta(x)\le\mathcal B(G(x))\le\mathcal B(B)=B.
\end{equation}
Thus $f(x)-B$ is a valid lower bound on the optimum.  The same
statement holds when a supplied diameter restricts the envelope,
because $\mathcal B_D(B)\le\mathcal B(B)$.
For a source with $\mathcal S(B)>0$, use a general valid certificate
such as the outer enclosure from Proposition~\ref{boundarygap:budget}
or the elementary bound in Corollary~\ref{boundarygap:elementary}.
An interval implementation may take the shortcut only after it has
certified the nonpositive sign of $\mathcal S(B)$; an unresolved sign
uses the general procedure.  The algebraic test still has a cost,
but it uses no objective, gradient or LMO query.  These statements
do not turn comparison counts into bit-complexity bounds.

\paragraph{The final phases of the adaptive work bound.}
If $g_{\rm c}>0$, strict increase of $\mathcal B$ and the theorem
give $\mathcal B^{-1}(u)=u$ for $0\le u\le g_{\rm c}$.
Thus the progress function in Theorem~\ref{nloracle:main}, with
$\Phi=\mathcal B$, becomes
\[
 \psi(u)=\kappa_\theta u
 \quad(0\le u\le g_{\rm c}),\qquad
 \kappa_\theta=\frac{1-\theta}{1+\theta}.
\]
With a multiplicative enclosure factor $\vartheta\ge1$, the
effective calibration $\Phi=\vartheta\mathcal B$ instead gives
\[
 \psi(u)=\kappa_\theta u/\vartheta
 \quad(0\le u\le\vartheta g_{\rm c}).
\]
Consequently the phases in that theorem whose upper error scale
$b_j$ lies in this range use these linear progress formulas and
$K(u)=H_0+H_1u$.  In particular,
$s_j=\kappa_\theta b_j/(2\vartheta)$ in the enclosed version.
No additional global gap-domination coefficient is needed for these
phases.  This is a consequence for the proved work bound, not a new
empirical convergence claim.
